\chapter{Diseases of the Blood and Hematopoietic System}
\label{ch:blood}
The hematopoietic system produces blood cells (red blood cells, white blood cells, platelets) in the bone marrow and includes the peripheral blood and spleen. Diseases of the blood include anemias, hemoglobinopathies, hemostatic disorders, and hematologic malignancies. These conditions are among the most prevalent diseases worldwide, affecting billions of individuals.

%-------------------------------------------------------------
\section{Anemias}
%-------------------------------------------------------------

%-------------------------------------------------------------

%-------------------------------------------------------------

%-------------------------------------------------------------

Anemias comprise a diverse group of hematological disorders characterized by a reduction in red blood cell mass, hemoglobin concentration, or hematocrit below normal reference limits. This deficiency impairs systemic oxygen delivery to peripheral tissues, precipitating tissue hypoxia and compensatory hemodynamic responses. Etiologies are broadly categorized into impaired erythrocyte production, accelerated erythrocyte destruction (hemolysis), and acute or chronic blood loss.

%-------------------------------------------------------------

%-------------------------------------------------------------

%-------------------------------------------------------------

%-------------------------------------------------------------

\label{sec:Anemias}
%-------------------------------------------------------------%-------------------------------------------------------------

\subsection{Aplastic Anemia}
\label{sec:Aplastic Anemia}
Aplastic anemia is a life-threatening bone marrow failure syndrome characterized by pancytopenia (anemia, leukopenia, thrombocytopenia) with a hypocellular bone marrow. The condition may be acquired (most common), inherited (Fanconi anemia, dyskeratosis congenita), or idiopathic, with acquired causes including autoimmune (T-cell mediated attack on hematopoietic stem cells), infections (hepatitis-associated aplastic anemia, parvovirus B19), drugs and toxins (chloramphenicol, chemotherapy agents, benzene), radiation, and pregnancy. Clinical features include fatigue, pallor (anemia), recurrent infections (neutropenia), and bleeding (thrombocytopenia). Diagnosis requires bone marrow biopsy showing markedly hypocellular marrow with increased fat cells, with severity classified as severe (marrow cellularity $<$25\%, or 25--50\% with $<$30\% residual hematopoiesis and two of three cytopenias: ANC $<$500, platelets $<$20,000, reticulocytes $<$1\%) or very severe (ANC $<$200). Management options include allogeneic hematopoietic stem cell transplantation (curative, preferred for young patients with matched donors) and immunosuppressive therapy (antithymocyte globulin + cyclosporine, effective for older patients), with eltrombopag (thrombopoietin receptor agonist) showing efficacy in combination with immunosuppressive therapy. Prognosis is guarded without treatment but improved with stem cell transplantation.

\subsection{Folate Deficiency Anemia}
\label{sec:Folate Deficiency Anemia}
Folate deficiency anemia is a megaloblastic anemia that is clinically indistinguishable from B12 deficiency anemia but without the neurological manifestations. The condition results from folate (vitamin B9) deficiency, with causes including poor dietary intake (alcoholism, inadequate fruit/vegetable consumption), malabsorption (celiac disease, tropical sprue), increased demand (pregnancy, hemolytic anemias), and medications (methotrexate, trimethoprim, phenytoin). Clinical features include fatigue, pallor, and glossitis. Diagnosis is by macrocytic anemia, low serum folate, and elevated homocysteine (with normal methylmalonic acid, distinguishing it from B12 deficiency). Management involves oral folic acid supplementation (1 mg daily) and addressing the underlying cause. Prognosis is excellent with treatment.

\subsection{Hemolytic Anemias}
\label{sec:Hemolytic Anemias}
Hemolytic anemias are a group of disorders characterized by premature destruction of red blood cells (RBCs), leading to anemia, elevated LDH, indirect hyperbilirubinemia (unconjugated), and reticulocytosis. The condition is classified as intrinsic (membrane disorders: hereditary spherocytosis, hereditary elliptocytosis; enzyme disorders: G6PD deficiency, pyruvate kinase deficiency; hemoglobinopathies: sickle cell disease, thalassemia) or extrinsic (immune: autoimmune hemolytic anemia, hemolytic disease of the newborn; mechanical: microangiopathic hemolytic anemia in TTP/HUS, DIC, prosthetic valves; infectious: malaria, babesiosis). Clinical features include fatigue, jaundice, dark urine, and splenomegaly. Diagnosis involves reticulocyte count, peripheral blood smear, direct antiglobulin test (Coombs test), hemoglobin electrophoresis, and specific tests based on suspected etiology. Management varies by cause and may include transfusions, corticosteroids (for immune hemolysis), splenectomy (for spherocytosis and some other hemolytic anemias), and disease-specific therapies. Prognosis depends on the underlying etiology.

\subsection{Iron Deficiency Anemia}
\label{sec:Iron Deficiency Anemia}
Iron deficiency anemia is a nutritional deficiency disorder that affects approximately 1.2 billion people worldwide, making it the most common nutritional deficiency globally. The condition results from insufficient iron for hemoglobin synthesis, leading to microcytic, hypochromic anemia, with causes including chronic blood loss (menstruation, GI bleeding from peptic ulcers, colorectal cancer, or NSAID use), inadequate dietary intake (vegetarians/vegans), and increased demand (pregnancy, growth). Clinical features include fatigue, pallor, dyspnea on exertion, pica (craving for non-food items like ice or clay), koilonychia (spoon nails), angular cheilitis, and glossitis. Diagnosis is by CBC showing microcytic anemia (low MCV, MCH, MCHC), low serum ferritin (most specific for iron deficiency), low serum iron, elevated TIBC, and low transferrin saturation, with peripheral smear showing hypochromic microcytes, pencil cells, and target cells. Management involves identifying and treating the underlying cause, and iron supplementation (oral ferrous sulfate 325 mg daily, taken with vitamin C to enhance absorption), with IV iron used for intolerance, malabsorption, or chronic kidney disease. Prognosis is excellent with appropriate treatment.

\subsection{Sideroblastic Anemia}
\label{sec:Sideroblastic Anemia}
Sideroblastic anemia is a disorder characterized by defective iron utilization in erythropoiesis, leading to iron accumulation in mitochondria of erythroid precursors and forming ring sideroblasts in the bone marrow. The condition may be hereditary (X-linked sideroblastic anemia, ALAS2 mutation) or acquired (myelodysplastic syndrome, alcohol, lead poisoning, isoniazid, pyridoxine deficiency). Clinical features include symptoms of anemia. Diagnosis shows a dimorphic picture (microcytic and macrocytic cells on peripheral smear), basophilic stippling, and a characteristic ring sideroblast pattern on Prussian blue staining of bone marrow. Management includes pyridoxine (vitamin B6, especially for hereditary forms), blood transfusions, and iron chelation therapy for iron overload, with allogeneic stem cell transplantation potentially curative. Prognosis depends on the underlying cause and response to treatment.

\subsection{Vitamin B12 Deficiency Anemia}
\label{sec:Vitamin B12 Deficiency Anemia}
Vitamin B12 deficiency anemia is a megaloblastic anemia that can cause irreversible neurological damage. The condition results from vitamin B12 (cobalamin) deficiency, with causes including pernicious anemia (autoimmune destruction of gastric parietal cells producing intrinsic factor), malabsorption (gastrectomy, ileal resection, Crohn's disease, bacterial overgrowth, celiac disease), dietary deficiency (strict vegans), and medications (metformin, proton pump inhibitors). Clinical features include fatigue, weakness, jaundice (mild, from ineffective erythropoiesis), glossitis (beefy red tongue), and neurological manifestations (subacute combined degeneration of the spinal cord---demyelination of dorsal columns and lateral corticospinal tracts, peripheral neuropathy, cognitive impairment, depression). Diagnosis is by macrocytic anemia (MCV $>$100 fL), hypersegmented neutrophils on peripheral smear, low serum B12, elevated methylmalonic acid and homocysteine, and elevated LDH. Management involves intramuscular cyanocobalamin or hydroxocobalamin injections, initially daily for one week, then monthly lifelong for pernicious anemia. Prognosis is good with treatment, though neurological damage may be partially irreversible if prolonged.

%-------------------------------------------------------------
\section{Coagulation Disorders}
%-------------------------------------------------------------

%-------------------------------------------------------------

%-------------------------------------------------------------

%-------------------------------------------------------------

%-------------------------------------------------------------

%-------------------------------------------------------------

%-------------------------------------------------------------

%-------------------------------------------------------------

\label{sec:Coagulation Disorders}
%-------------------------------------------------------------%-------------------------------------------------------------

\subsection{Disseminated Intravascular Coagulation}
\label{sec:Disseminated Intravascular Coagulation}
Disseminated intravascular coagulation is an acquired, life-threatening coagulation disorder characterized by systemic activation of coagulation, leading to widespread microvascular thrombosis and consumption of platelets and clotting factors, resulting in a paradoxical bleeding tendency. The condition is always secondary to an underlying condition: sepsis (most common cause), trauma, malignancy, obstetric complications (placental abruption, amniotic fluid embolism), transfusion reactions, and severe liver disease. Clinical features include widespread bleeding (from IV sites, surgical wounds, mucosal surfaces), microvascular thrombosis (organ dysfunction), and hemodynamic instability. Diagnosis is based on clinical findings and laboratory evidence: thrombocytopenia, prolonged PT and aPTT, low fibrinogen, elevated D-dimer and fibrin degradation products, and schistocytes on peripheral smear, with the ISTH DIC score aiding diagnosis. Management targets the underlying cause, with supportive care including platelet and FFP transfusions for active bleeding, cryoprecipitate for fibrinogen $<$100 mg/dL, and heparin for thrombotic predominant DIC. Prognosis depends on the underlying cause and severity.

\subsection{Hemophilia A}
\label{sec:Hemophilia A}
Hemophilia A is an X-linked recessive bleeding disorder caused by deficiency or dysfunction of coagulation factor VIII, affecting approximately 1 in 5000 males. The condition results from factor VIII deficiency, with severity classified by residual factor VIII activity: severe ($<$1\%), moderate (1--5\%), mild (5--40\%). Clinical features include spontaneous hemarthroses (knee, elbow, ankle---leading to chronic hemophilic arthropathy), deep muscle hematomas, retroperitoneal hemorrhage, and prolonged bleeding after surgery or trauma, with intracranial hemorrhage being the leading cause of death. Diagnosis is by prolonged aPTT with normal PT, decreased factor VIII activity, and positive family history. Management involves factor VIII replacement therapy (recombinant or plasma-derived), given prophylactically (for severe hemophilia, typically 3 times weekly) or on-demand for bleeding episodes, with emicizumab (a bispecific antibody bridging factor IXa and factor X) used for prophylaxis. Prognosis has been transformed by factor replacement and gene therapy (valoctocogene roxaparvovec), though inhibitors (neutralizing alloantibodies against factor VIII) develop in 25--30\% of severe hemophilia A and require bypassing agents or immune tolerance induction.

\subsection{Hemophilia B}
\label{sec:Hemophilia B}
Hemophilia B (Christmas disease) is an X-linked recessive bleeding disorder caused by deficiency or dysfunction of coagulation factor IX, accounting for approximately 15\% of hemophilia cases. Clinical features and classification are identical to hemophilia A. Diagnosis is by prolonged aPTT, normal PT, and decreased factor IX activity. Management involves factor IX replacement therapy (recombinant or plasma-derived), with a longer half-life than factor VIII allowing less frequent dosing, and emicizumab is not effective for hemophilia B. Prognosis has improved with gene therapy (etranacogene dezaparvovec), which offers durable factor IX expression and reduction in bleeding episodes.

\subsection{Immune Thrombocytopenic Purpura}
\label{sec:Immune Thrombocytopenic Purpura}
Immune thrombocytopenic purpura is an autoimmune disorder characterized by isolated thrombocytopenia (platelet count $<$100,000/$\mu$L) due to antibody-mediated platelet destruction and impaired megakaryopoiesis. The condition may be primary (idiopathic) or secondary (associated with SLE, antiphospholipid syndrome, infections [HIV, H.~pylori], drugs, lymphoproliferative disorders, or vaccination). Clinical features include petechiae, purpura, mucosal bleeding (epistaxis, gingival bleeding, menorrhagia), and in severe cases, intracranial hemorrhage (rare but life-threatening). Diagnosis is one of exclusion: isolated thrombocytopenia on CBC with normal WBC and hemoglobin, normal peripheral smear (except for thrombocytopenia), and exclusion of other causes. First-line management includes corticosteroids (prednisone or dexamethasone) and IVIG or anti-D immunoglobulin for rapid platelet increase, with second-line therapies including thrombopoietin receptor agonists (romiplostim, eltrombopag, avatrombopag), rituximab, splenectomy, and fostamatinib (SYK inhibitor). Prognosis is generally good, with TPO receptor agonists having largely replaced splenectomy as the preferred second-line therapy.

\subsection{Thrombotic Thrombocytopenic Purpura}
\label{sec:Thrombotic Thrombocytopenic Purpura}
Thrombotic thrombocytopenic purpura (TTP) is a rare, life-threatening thrombotic microangiopathy characterized by microvascular thrombosis, thrombocytopenia, and microangiopathic hemolytic anemia, with an incidence of 3--4 per million. The condition results from severe deficiency of ADAMTS13, a von Willebrand factor (VWF) cleaving protease: immune-mediated (acquired TTP, IgG autoantibodies, adulthood) or congenital (Upshaw-Schulman syndrome, ADAMTS13 gene mutations, childhood). ADAMTS13 deficiency leads to accumulation of ultralarge VWF multimers, causing spontaneous platelet adhesion and microthrombi formation in arterioles and capillaries. Clinical features include the classic pentad: thrombocytopenia (petechiae, purpura, bleeding), microangiopathic hemolytic anemia (fever, pallor, jaundice), neurologic involvement (headache, confusion, seizures, focal deficits, coma), renal involvement (elevated creatinine, hematuria), and fever. Not all five features are present simultaneously; thrombocytopenia and MAHA are sufficient to suspect TTP. Diagnosis is based on peripheral blood smear showing schistocytes (helmet cells, fragmented RBCs), elevated LDH, low haptoglobin, negative direct Coombs test, and ADAMTS13 activity <10\%. Management is a medical emergency: urgent plasma exchange (PLEX, 1--1.5 plasma volume daily) is the cornerstone of therapy, combined with high-dose corticosteroids (prednisone 1 mg/kg). For refractory/relapsed immune TTP: rituximab, caplacizumab (anti-VWF nanobody, blocks platelet adhesion), or immunosuppression (vincristine, cyclophosphamide). Prognosis has improved dramatically: mortality declined from >90\% to 10--20\% with prompt PLEX.

\subsection{Von Willebrand Disease}
\label{sec:Von Willebrand Disease}
Von Willebrand disease is the most common inherited bleeding disorder, affecting up to 1\% of the population. The condition results from quantitative or qualitative defects in von Willebrand factor (VWF), a large multimeric glycoprotein essential for primary hemostasis (platelet adhesion to injured vessel walls) and carrier/stabilizer of factor VIII, with type 1 (partial quantitative deficiency, 70--80\%) being the most common and mildest, type 2 (qualitative deficiency) having various functional abnormalities, and type 3 (complete deficiency) being the most severe. Clinical features include mucocutaneous bleeding: epistaxis, menorrhagia, easy bruising, prolonged bleeding after dental procedures or surgery, and GI bleeding. Diagnosis is by decreased VWF antigen, decreased VWF activity (ristocetin cofactor assay), and decreased or normal factor VIII. Management includes desmopressin (DDAVP, for type 1 and some type 2A patients), VWF-containing factor VIII concentrates (for type 3 or DDAVP non-responders), and antifibrinolytics (tranexamic acid, for mucosal bleeding). Prognosis is generally good with appropriate management.

%-------------------------------------------------------------
\section{Hematologic Malignancies}
%-------------------------------------------------------------

%-------------------------------------------------------------

%-------------------------------------------------------------

%-------------------------------------------------------------

%-------------------------------------------------------------

%-------------------------------------------------------------

%-------------------------------------------------------------

%-------------------------------------------------------------

\label{sec:Hematologic Malignancies}
%-------------------------------------------------------------%-------------------------------------------------------------

\subsection{Acute Lymphoblastic Leukemia}
\label{sec:Acute Lymphoblastic Leukemia}
Acute lymphoblastic leukemia is the most common childhood cancer (peak age 2--5 years), characterized by proliferation of immature lymphoid precursors (blasts) in the bone marrow, with B-cell ALL accounting for approximately 85\% of childhood cases. Clinical features include bone marrow failure: anemia (fatigue, pallor), neutropenia (fever, infections), and thrombocytopenia (bleeding, bruising), along with bone pain, lymphadenopathy, hepatosplenomegaly, and testicular enlargement, with CNS involvement possibly presenting with headache, vomiting, or cranial nerve palsies. Diagnosis is by bone marrow biopsy showing $>$20\% blasts, immunophenotyping (flow cytometry), cytogenetics [hyperdiploidy, t(12;21) ETV6-RUNX1---favorable; t(9;22) BCR-ABL---unfavorable; MLL rearrangements, Philadelphia chromosome-like ALL], and molecular studies. Management is risk-stratified and prolonged (2--3 years): induction (vincristine, corticosteroids, L-asparaginase, $\pm$ anthracyclines), consolidation, interim maintenance, delayed intensification, and maintenance therapy, with CNS-directed therapy (intrathecal chemotherapy, $\pm$ cranial radiation) essential. Prognosis is excellent in children, with cure rates exceeding 90\% in developed countries, though adult ALL has a worse prognosis (40--50\% long-term survival), and CAR-T cell therapy (tisagenlecleucel) and blinatumomab are used for relapsed/refractory B-cell ALL.

\subsection{Acute Myeloid Leukemia}
\label{sec:Acute Myeloid Leukemia}
Acute myeloid leukemia is the most common acute leukemia in adults (median age 68), characterized by clonal proliferation of immature myeloid precursors (blasts) in the bone marrow, with risk factors including prior chemotherapy (alkylating agents, topoisomerase II inhibitors), radiation, myelodysplastic syndrome, and myeloproliferative neoplasms. Clinical features include bone marrow failure (anemia, neutropenia, thrombocytopenia) and symptoms of organ infiltration (hepatosplenomegaly, gum hyperplasia, skin infiltration). Diagnosis requires bone marrow biopsy showing $\geq$20\% blasts, immunophenotyping, cytogenetics (t(15;17) APL---favorable; inv(16), t(8;21)---favorable; complex karyotype, monosomy 5/7---unfavorable), and molecular studies (FLT3, NPM1, CEBPA, IDH1/2, TP53 mutations). Management involves induction chemotherapy (cytarabine + anthracycline [``7+3''] or CPX-351 for secondary AML) followed by consolidation (cytarabine, allogeneic stem cell transplantation for intermediate/high-risk), with acute promyelocytic leukemia (APL, t(15;17), PML-RARA) treated with all-trans retinoic acid (ATRA) + arsenic trioxide with excellent cure rates ($>$90\%). Prognosis is variable depending on cytogenetic and molecular risk factors.

\subsection{Chronic Lymphocytic Leukemia}
\label{sec:Chronic Lymphocytic Leukemia}
Chronic lymphocytic leukemia is the most common leukemia in adults in Western countries, with a median age at diagnosis of 70 years, characterized by accumulation of mature-appearing but functionally incompetent B lymphocytes in the blood, bone marrow, lymph nodes, and spleen. Clinical features include generalized lymphadenopathy, hepatosplenomegaly, fatigue, weight loss, night sweats, and autoimmune complications (autoimmune hemolytic anemia, immune thrombocytopenia), with many patients being asymptomatic and diagnosed incidentally on routine CBC showing lymphocytosis ($>$5000/$\mu$L B lymphocytes). Diagnosis is by flow cytometry showing clonal B cells co-expressing CD5, CD19, CD20 (dim), CD23, with prognostic markers including IGHV mutation status (mutated = favorable, unmutated = unfavorable), TP53 mutations/deletions (very unfavorable), del(13q) (favorable), trisomy 12, and del(11q). Management is indicated for symptomatic or bulky disease, progressive cytopenias, or autoimmune complications, with options including BTK inhibitors (ibrutinib, acalabrutinib, zanubrutinib), BCL-2 inhibitor (venetoclax) plus obinutuzumab, and FCR (fludarabine, cyclophosphamide, rituximab) for young patients with mutated IGHV. Prognosis is generally favorable, with CLL being incurable but highly treatable and many patients having prolonged survival.

\subsection{Essential Thrombocythemia}
\label{sec:Essential Thrombocythemia}
Essential thrombocythemia is a myeloproliferative neoplasm characterized by sustained thrombocytosis (platelet count $>$450,000/$\mu$L) with increased megakaryocyte proliferation, in the absence of other myeloproliferative neoplasms. The condition is driven by JAK2 V617F (50\%), CALR (25--30\%), or MPL (3--5\%) mutations; approximately 10\% are triple-negative. Clinical features may include microvascular symptoms (erythromelalgia, headache, visual disturbances), thrombosis (arterial or venous), or hemorrhage (with very high platelet counts $>$1,000,000/$\mu$L due to acquired von Willebrand syndrome), though many patients are asymptomatic. Diagnosis is based on sustained thrombocytosis with the characteristic mutation profile. Management of low-risk ET involves observation with low-dose aspirin; high-risk ET (age $>$60 or history of thrombosis) requires cytoreduction with hydroxyurea or anagrelide. Prognosis is favorable with appropriate management.

\subsection{Hodgkin Lymphoma}
\label{sec:Hodgkin Lymphoma}
Hodgkin lymphoma is a malignant lymphoid neoplasm characterized by Reed-Sternberg cells (large, binucleated or multinucleated cells with prominent nucleoli) in a background of reactive inflammatory cells, with two major types: classical HL (nodular sclerosis, mixed cellularity, lymphocyte-rich, lymphocyte-depleted---accounting for 95\% of cases) and nodular lymphocyte-predominant HL (NLPHL), with a bimodal age distribution (peaks at 20--30 and $>$60 years). Clinical features typically include painless lymphadenopathy (cervical, supraclavicular, mediastinal), B symptoms (fever $>$38$^\circ$C, drenching night sweats, weight loss $>$10\% in 6 months), and pruritus. Diagnosis requires lymph node biopsy showing Reed-Sternberg cells (CD15+, CD30+, CD45-), with PET-CT used for staging (Lugano classification). Management of early-stage favorable disease involves combined modality therapy (ABVD chemotherapy [doxorubicin, bleomycin, vinblastine, dacarbazine] + involved-field radiation therapy), while advanced-stage disease uses ABVD or escalated BEACOPP with PET-adapted therapy. Prognosis is excellent, with cure rates exceeding 80\%; brentuximab vedotin and checkpoint inhibitors (nivolumab, pembrolizumab) are used for relapsed/refractory disease.

\subsection{Multiple Myeloma}
\label{sec:Multiple Myeloma}
Multiple myeloma is a plasma cell neoplasm that accounts for 1\% of all cancers and 13\% of hematologic malignancies, with a median age at diagnosis of 70 years. The condition results from clonal proliferation of malignant plasma cells in the bone marrow, production of monoclonal immunoglobulin (M-protein), and end-organ damage summarized by the CRAB criteria: Calcium elevation (hypercalcemia), Renal insufficiency, Anemia, and Bone lesions (lytic lesions on imaging). Clinical features include bone pain (vertebral compression fractures, pathologic fractures), fatigue (anemia), recurrent infections (immunoparesis---suppression of normal immunoglobulins), and renal failure. Diagnosis requires $\geq$10\% clonal bone marrow plasma cells plus evidence of end-organ damage (CRAB) or biomarkers of malignancy (bone marrow plasma cells $\geq$60\%, free light chain ratio $\geq$100, or $>$1 focal lesion on MRI). Risk stratification uses the International Staging System (ISS) and revised ISS (R-ISS), with del(17p), t(4;14), and t(14;16) as high-risk cytogenetic abnormalities. Management is risk-adapted and typically involves induction (bortezomib, lenalidomide, dexamethasone---VRd), followed by autologous stem cell transplantation for eligible patients, followed by maintenance therapy (lenalidomide). Prognosis has improved to a median survival of 7--10 years with modern therapy, which includes novel agents such as proteasome inhibitors, immunomodulatory drugs, monoclonal antibodies (daratumumab), bispecific antibodies, and CAR-T cell therapy.

\subsection{Myelodysplastic Syndromes}
\label{sec:Myelodysplastic Syndromes}
Myelodysplastic syndromes are a heterogeneous group of clonal hematopoietic stem cell disorders that predominantly affect the elderly (median age 70--75). The condition is characterized by dysplastic and ineffective hematopoiesis, peripheral cytopenias, and a risk of transformation to acute myeloid leukemia (AML). Clinical features include symptoms of cytopenias: fatigue (anemia), infections (neutropenia), and bleeding (thrombocytopenia). Diagnosis requires bone marrow biopsy showing dysplasia ($\geq$10\% of cells in one or more myeloid lineages), blast percentage ($<$20\%---above 20\% classifies as AML), and cytogenetic/molecular abnormalities (del(5q), del(7q), del(20q), SF3B1, TET2, ASXL1 mutations). Risk stratification uses the International Prognostic Scoring System (IPSS-R, now IPSS-M). Management of low-risk MDS includes supportive care, erythropoiesis-stimulating agents ($\pm$ lenalidomide for del(5q) MDS), and luspatercept, while higher-risk MDS is treated with hypomethylating agents (azacitidine, decitabine). Prognosis depends on risk category, with allogeneic stem cell transplantation being the only curative option, typically for patients $<$75 with a donor.

\subsection{Myelofibrosis}
\label{sec:Myelofibrosis}
Myelofibrosis is a myeloproliferative neoplasm characterized by clonal megakaryocyte proliferation, bone marrow fibrosis (reticulin and/or collagen deposition), extramedullary hematopoiesis (splenomegaly, hepatomegaly), and progressive cytopenias, which may arise de novo (primary MF) or evolve from PV or ET (post-PV/post-ET MF). The condition is driven by JAK2, CALR, or MPL mutations in over 90\% of cases, with ASXL1, EZH2, SRSF2 mutations associated with worse prognosis. Clinical features include massive splenomegaly (early satiety, abdominal discomfort), constitutional symptoms (fatigue, weight loss, night sweats, fever), bone pain, and progressive anemia requiring transfusions. Diagnosis requires bone marrow biopsy showing megakaryocyte proliferation and atypia, reticulin/collagen fibrosis, and exclusion of other causes. The Dynamic International Prognostic Scoring System (DIPSS) and molecular-enhanced versions (MIPSS70) guide treatment. Management options include ruxolitinib (JAK1/2 inhibitor, improves splenomegaly and symptoms), fedratinib (JAK2 inhibitor), momelotinib (JAK1/2 inhibitor with anemia benefit), and allogeneic stem cell transplantation (the only curative option). Prognosis is variable depending on risk category.

\subsection{Non-Hodgkin Lymphoma}
\label{sec:Non-Hodgkin Lymphoma}
Non-Hodgkin lymphoma encompasses a heterogeneous group of over 60 subtypes of lymphoid malignancies not classified as Hodgkin lymphoma, with the most common subtypes including diffuse large B-cell lymphoma (DLBCL, most common NHL overall, aggressive), follicular lymphoma (most common indolent NHL), and mantle cell lymphoma, with risk factors including immunosuppression (HIV/AIDS---Burkitt lymphoma, EBV-associated lymphomas), autoimmune diseases, and \textit{H.~pylori} infection (gastric MALT lymphoma). Clinical features vary by subtype: aggressive lymphomas (DLBCL) present with rapidly growing masses, B symptoms, and organ compromise; indolent lymphomas (follicular) present with waxing and waning painless lymphadenopathy. Diagnosis is by lymph node biopsy with histopathology, immunophenotyping, cytogenetics, and molecular studies. Management of DLBCL is R-CHOP (rituximab, cyclophosphamide, doxorubicin, vincristine, prednisone), with approximately 60--70\% cure rate, while indolent NHL may be observed (watch and wait) or treated with rituximab-based regimens. Prognosis varies widely by subtype.

\subsection{Polycythemia Vera}
\label{sec:Polycythemia Vera}
Polycythemia vera is a Philadelphia chromosome-negative myeloproliferative neoplasm characterized by clonal proliferation of hematopoietic stem cells, predominantly producing excess red blood cells, and often accompanied by leukocytosis and thrombocytosis, with over 95\% of cases harboring the JAK2 V617F mutation (or JAK2 exon 12 mutations). Clinical features include erythromelalgia (burning pain and redness of extremities), pruritus (especially aquagenic---after bathing), headache, plethora, hepatosplenomegaly, and thrombotic or hemorrhagic complications. The WHO 2022 diagnostic criteria require hemoglobin $>$16.5 g/dL in men or $>$16.0 g/dL in women, or increased red cell mass, plus bone marrow biopsy showing hypercellularity with trilineage growth, or JAK2 mutation. Management includes phlebotomy (target hematocrit $<$45\%), low-dose aspirin, and cytoreductive therapy (hydroxyurea for high-risk patients), with ruxolitinib (JAK1/2 inhibitor) used for hydroxyurea-resistant or intolerant disease. Prognosis is generally favorable but carries a risk of transformation to myelofibrosis (10--20\%) or AML (5--10\%).

%-------------------------------------------------------------
\section{Hemoglobinopathies}
%-------------------------------------------------------------

%-------------------------------------------------------------

%-------------------------------------------------------------

%-------------------------------------------------------------

Hemoglobinopathies are genetic disorders affecting the synthesis, structure, or production rate of hemoglobin molecules within erythrocytes. Structural hemoglobin variants result from amino acid substitutions in globin chains, whereas thalassemias arise from quantitative reduction or absent synthesis of specific globin polypeptide chains. These mutations impair oxygen transport dynamics, reduce red cell lifespan, and frequently cause chronic hemolytic anemia and microvascular occlusion.

%-------------------------------------------------------------

%-------------------------------------------------------------

%-------------------------------------------------------------

%-------------------------------------------------------------

\label{sec:Hemoglobinopathies}
%-------------------------------------------------------------%-------------------------------------------------------------

\subsection{Hereditary Spherocytosis}
\label{sec:Hereditary Spherocytosis}
Hereditary spherocytosis is the most common inherited hemolytic anemia in individuals of Northern European descent, affecting approximately 1 in 2000. The condition results from mutations in genes encoding red blood cell membrane proteins (spectrin, ankyrin, band 3, protein 4.2), leading to defective vertical interactions between the membrane and cytoskeleton, producing spherocytic, less deformable RBCs that are trapped and destroyed in the spleen. Clinical features range from asymptomatic (compensated hemolysis) to severe hemolytic anemia with jaundice, splenomegaly, gallstones, and aplastic crisis (parvovirus B19 infection). Diagnosis is by osmotic fragility testing (increased fragility), EMA (eosin-5-maleimide) binding test by flow cytometry (reduced fluorescence), and peripheral smear showing spherocytes. Management ranges from observation and folic acid supplementation to splenectomy (for moderate-to-severe cases, ideally deferred until age 5 to reduce OPSI risk). Prognosis is good with appropriate management.

\subsection{Methemoglobinemia}
\label{sec:Methemoglobinemia}
Methemoglobinemia is a disorder characterized by increased levels of methemoglobin (ferric Hb, Fe$^{3+}$), which is unable to bind oxygen and causes functional anemia, with a normal methemoglobin level <1\% of total hemoglobin. The condition results from either congenital deficiency of cytochrome b5 reductase (NADH-methemoglobin reductase, autosomal recessive, most common congenital cause) or exposure to oxidizing agents (nitrates, nitrites, dapsone, benzocaine, sulfonamides, aniline dyes). Clinical features correlate with methemoglobin levels: 10--20\% causes cyanosis without symptoms; 20--50\% causes headache, dyspnea, tachycardia, and altered mental status; >50\% causes seizures, coma, cardiac arrhythmias, and death. A characteristic finding is chocolate-brown discoloration of blood that does not turn red when exposed to air, and pulse oximetry showing a saturation of approximately 85\% that does not change with oxygen supplementation. Diagnosis is confirmed by co-oximetry (multi-wavelength blood gas analysis) showing elevated methemoglobin. Management includes supplemental oxygen, intravenous methylene blue (1--2 mg/kg over 5 minutes) for symptomatic or severe methemoglobinemia (>20--30\%), which acts as an electron donor to reduce methemoglobin to hemoglobin. Methylene blue ineffective in G6PD deficiency (risk of hemolysis); alternative is ascorbic acid or exchange transfusion. Prognosis is excellent with treatment; untreated severe methemoglobinemia can be fatal.

\subsection{Sickle Cell Disease}
\label{sec:Sickle Cell Disease}
Sickle cell disease is a group of inherited hemoglobinopathies predominantly affecting individuals of African, Mediterranean, Middle Eastern, and Indian descent. The condition results from mutations in the $\beta$-globin gene, leading to production of abnormal hemoglobin S (HbS), with the most common genotype being HbSS (sickle cell anemia, caused by homozygosity for the Glu6Val mutation); deoxygenated HbS polymerizes, causing red blood cells to assume a sickle shape, leading to vaso-occlusion, hemolysis, and end-organ damage. Clinical features include vaso-occlusive crises (acute pain crises, the hallmark of SCD), acute chest syndrome (chest pain, fever, pulmonary infiltrates---a leading cause of hospitalization and death), stroke (both ischemic and hemorrhagic, occurring in up to 20\% of children), avascular necrosis of the femoral head, priapism, splenic sequestration and functional asplenia, delayed growth, and chronic kidney disease. Diagnosis is by hemoglobin electrophoresis and HPLC showing elevated HbS. Management includes hydroxyurea (increases fetal hemoglobin [HbF], reduces crisis frequency), chronic transfusion therapy, pain management for crises, and prophylactic penicillin (for children to prevent \textit{S.~pneumoniae} sepsis), with novel therapies including voxelotor (HbS polymerization inhibitor), crizanlizumab (anti-P-selectin antibody), and L-glutamine (oxidative stress reducer). Prognosis has improved significantly with modern therapies, with allogeneic stem cell transplantation offering cure but limited by donor availability, and gene therapy (lovotibeglogene autotemcel, exagamglogene autotemcel) showing remarkable efficacy in clinical trials.

\subsection{Thalassemia}
\label{sec:Thalassemia}
Thalassemia is a group of inherited hemoglobin disorders characterized by reduced or absent synthesis of $\alpha$- or $\beta$-globin chains, leading to ineffective erythropoiesis and hemolytic anemia. $\beta$-Thalassemia major (Cooley's anemia, homozygous $\beta^0$ mutations) presents in the first year of life with severe anemia, failure to thrive, hepatosplenomegaly, and skeletal deformities (``hair-on-end'' appearance on skull X-ray from marrow expansion). $\beta$-Thalassemia intermedia has milder anemia and may not require regular transfusions. $\alpha$-Thalassemia major (Hb Bart's hydrops fetalis, deletion of all four $\alpha$-globin genes) is uniformly fatal in utero. $\alpha$-Thalassemia trait (deletion of two genes) causes mild microcytic anemia. Diagnosis is by CBC, hemoglobin electrophoresis, and genetic testing. Management of $\beta$-thalassemia major requires lifelong transfusion therapy and iron chelation (deferasirox, deferoxamine, deferiprone) to prevent transfusional iron overload, which causes cardiomyopathy, liver cirrhosis, and endocrinopathies. Prognosis has improved with stem cell transplantation (curative for $\beta$-thalassemia major) and gene therapy (betibeglogene autotemcel).

%-------------------------------------------------------------
\section{Thrombotic Disorders}
%-------------------------------------------------------------

%-------------------------------------------------------------

%-------------------------------------------------------------

%-------------------------------------------------------------

%-------------------------------------------------------------

%-------------------------------------------------------------

%-------------------------------------------------------------

%-------------------------------------------------------------

\label{sec:Thrombotic Disorders}
%-------------------------------------------------------------%-------------------------------------------------------------

\subsection{Deep Vein Thrombosis}
\label{sec:Deep Vein Thrombosis}
Deep vein thrombosis is the formation of a blood clot (thrombus) in a deep vein, most commonly in the lower extremities (iliac, femoral, popliteal veins), affecting approximately 1 in 1000 people annually and being a major cause of morbidity and mortality due to its most feared complication: pulmonary embolism. Risk factors are described by Virchow's triad: venous stasis (immobilization, prolonged travel, paralysis, heart failure), endothelial injury (surgery, trauma, indwelling catheters), and hypercoagulability (malignancy, oral contraceptives, pregnancy, inherited thrombophilia [Factor V Leiden, prothrombin gene mutation], antiphospholipid syndrome). Clinical features include unilateral leg swelling, pain (often described as a dull ache or cramping), warmth, erythema, and palpable cord along the affected vein, though many DVTs are asymptomatic. Diagnosis involves clinical probability assessment (Wells score), D-dimer testing (high negative predictive value), and compression ultrasonography (first-line imaging, showing non-compressibility of the affected vein). Management is anticoagulation: initial therapy with low-molecular-weight heparin (enoxaparin), fondaparinux, or direct oral anticoagulants (DOACs---rivaroxaban, apixaban), followed by long-term anticoagulation for 3--6 months (first unprovoked DVT) or longer depending on risk of recurrence. Prognosis is good with treatment, though complications include post-thrombotic syndrome (chronic venous insufficiency in 20--50\% of patients) and chronic thromboembolic pulmonary hypertension.

